% ═══════════════════════════════════════════════════════════════
% Section 9: Limitations and Future Work
% ═══════════════════════════════════════════════════════════════

\section{Limitations and Future Work}
\label{sec:limitations}

We present these limitations transparently to guide future validation efforts and
to prevent overstatement of the framework's current maturity. QIF, TARA and NISS are
proposed and unvalidated: not independently peer-reviewed as a whole, and not adopted by
any standards body or device maker.

\subsection{No Empirical Validation on Real BCI Devices}

The \qif framework has not been validated against operational BCI hardware.
The TARA taxonomy was developed through literature review, threat modeling, and
systematic analysis rather than penetration testing of actual neural devices.
While the framework has been applied to one real software vulnerability
(Section~\ref{sec:case-studies}), this covers only device-side code. Validation
against interface and neural threats requires access to implanted BCI
patients and clinical environments---resources unavailable to independent
researchers.

\subsection{Reproducibility of the v\SnapshotRegistryVersion{} Figures}

For this revision the snapshot figures were recomputed from the registry file committed
with v1.5. The evidence-status, filed-severity, category, dual-use, consent-tier, PINS and
Neural Impact Chain counts reproduce exactly. Three do not: the gap groups (published
13/31/64, totalling \SnapshotGapTechniques{} or \SnapshotGapPercent\%; file 14/30/59,
totalling 103 or 94.5\%), the NISS severity distribution (published 19/37/52/1 high to none;
file 21/37/50/1), and several stored NISS scores, which do not reproduce from
Equation~\ref{eq:niss}. The cause was not identified. The text of this revision uses the
reproducible gap counts; the published ones are kept in the tables for the record. The case
studies and representative techniques of v1.5 did not match the registry at all and were
replaced (Sections~\ref{sec:tara} and~\ref{sec:case-studies}).

\subsection{Diagnostic Category Associations Not Clinically Validated}

The Neural Impact Chain associates techniques with diagnostic category references through
neuroanatomical addresses and functional neuroscience. These associations have not been
reviewed by psychiatrists, neurologists or clinical neuroscientists. They are references for
threat modeling, not diagnostic claims, and must not inform clinical decision-making.

\subsection{NISS Weights Not Calibrated}

The default profile weights BI, CV, RV and NP at 1.0 and CR and CD at 0.5
(Equation~\ref{eq:niss}). Neither these nor the four context profiles have been calibrated
against empirical data, expert elicitation or clinical outcomes, and the read-only ceiling
of 3.0 (Section~\ref{sec:niss-limits}) is a direct consequence of the default choice.
Calibration requires expert panel scoring of representative scenarios, sensitivity analysis
across weight configurations, and correlation with observed outcomes when available.

\subsection{No Interrater Reliability Study}

NISS scores in the TARA registry were assigned by a single analyst (the author) with AI
assistance. No interrater reliability study has been conducted. CVSS interrater reliability
is a known challenge~\cite{first2023cvss4}; NISS, with its novel neural-specific metrics,
likely faces greater variability. A formal study with domain experts from both
cybersecurity and neuroscience is needed.

\subsection{Evidence and Completeness}

The analysed snapshot contains \SnapshotTechniques{} techniques (v\SnapshotRegistryVersion{}),
of which 31 are theoretical and 1 speculative. The live registry holds \LiveTechniques{}
techniques and \LiveVariants{} variant across \LiveTactics{} tactics
(v\LiveRegistryVersion{}), of which 95 are theoretical. \LiveWithCvss{} carry CVSS v4.0 base
vectors; of the NISS scores \LiveNissScored{} are complete, \LiveNissProvisional{} are
provisional ranges and \LiveNissUnscored{} are unscored under the evidence rule. The
\LiveNissScored{} scored techniques were scored before that rule and have not been
re-examined against it, so the scored count is expected to fall. Live figures are not
substituted into the analysis, because every proportion in
Sections~\ref{sec:tara}--\ref{sec:governance} has a numerator computed over the snapshot;
re-pointing the denominator would produce figures that describe no actual registry state.
Current values are derivable from \texttt{datalake/qtara-registrar.json} in the project
repository. The addressing model has not been reviewed by a human neuroanatomist, and the
catalogue records techniques one person at a time, so attacks spread across a population are
outside its model (Section~\ref{sec:related-work}).

\subsection{Single-Author Bias}

The framework was developed by a single independent researcher. While multi-model AI
review was used throughout development (Claude, Gemini, ChatGPT), the architectural
decisions, scoring assignments, and clinical associations reflect a single perspective, and
AI review does not substitute for independent human review.

\subsection{AI Tool Disclosure}

In accordance with arXiv, ACM, and IEEE policies on AI-assisted research, we
disclose the following. Large language models (Claude~3.5/4, Gemini~2.0, ChatGPT-4o)
were used during the development of this framework and versions up to v1.5. Their roles, by
section:

\begin{itemize}
  \item \textbf{Literature review} (Sections~2, 3, 7): AI assisted in identifying
        relevant publications and summarizing prior work. All citations were
        independently verified by the author via DOI resolution and publisher page
        confirmation (see version history below).
  \item \textbf{Code generation} (figures, TARA data pipeline): AI generated
        Python scripts for data analysis, visualization, and figure generation. All
        code was reviewed and modified by the author.
  \item \textbf{Writing assistance} (all sections): AI provided draft text that was
        substantially rewritten, restructured, and verified by the author. For versions
        up to v1.5, the proportion of AI-generated text retained verbatim is estimated at
        less than 15\%.
  \item \textbf{Cross-validation}: AI models were used to cross-check factual claims
        across multiple systems (Claude, Gemini, ChatGPT) to reduce single-model
        bias. Disagreements between models were resolved by the author through
        primary source verification.
\end{itemize}

\noindent\textbf{This revision (v1.6).} The v1.6 changes were drafted on 2026-10-17 by Claude
Opus~5.5 (Anthropic), operating through Claude Code, from the project's data files, the v9
whitepaper and the repository history. The AI system wrote the new and rewritten text (the
abstract; Section~\ref{sec:why-neurosecurity}; the cognitive-neurosecurity subsection of
Section~\ref{sec:related-work}; Section~\ref{sec:hourglass}; the evidence, hazards and
representative-technique subsections of Section~\ref{sec:tara}; the gap analysis and
Section~\ref{sec:niss-limits}; the walkthrough and association subsections of
Section~\ref{sec:nic}; Sections~\ref{sec:why-neuroethics} and the remainder of
Section~\ref{sec:governance}; Section~\ref{sec:case-studies}; this section's reproducibility,
evidence and version notes; and Section~\ref{sec:conclusion}). It recomputed every figure
it added from the data files with scripts it ran, computed the CVSS scores in
Section~\ref{sec:case-studies}, and re-resolved each added citation through the Crossref or
arXiv API on the same date. Text new in v1.6 is predominantly AI-drafted and was pending
the author's review when drafted; it should not be read as author-reviewed until the
Zenodo version description says so.

\noindent\textbf{Human-originated contributions:} The original band model, the TARA
taxonomy structure, the NISS scoring methodology, the Neural Impact Chain method, all
diagnostic category associations, all NISS metric assignments, and the research conclusions
of v1.0--v1.5 were conceived, designed and decided by the author. The addressing revision,
the evidence rule and the retirement of the confirmed status were made in the project's data
before this revision, which reports them.

\noindent The author takes full responsibility for all content in this paper,
irrespective of how it was generated. A complete, auditable transparency log
documenting AI contributions and verification steps is maintained at
\url{https://github.com/qinnovates/qinnovate/blob/main/governance/TRANSPARENCY.md}.

\paragraph{Version history.}
An earlier version of this preprint (v1.0) contained citation errors
introduced during AI-assisted bibliography construction, including three
fabricated entries. These were corrected in v1.1 through a two-pass
independent verification audit. Version 1.2 corrected an internal percentage
inconsistency and added the author responsibility statement above.
Version~1.3 expanded the regulatory context with FDORA/PATCH Act Section~524B analysis and
added Schroder et al.~(2025) to the related work. Version~1.4 added the Sherman \&
Guillery~(2006) citation, replaced a mislabeled DBS citation with Hallett~(2007) for TMS,
corrected case study scores, fixed the dual-use classification label, and completed a
citation and claims verification pass. Version~1.5 expanded the registry from 102 to
\SnapshotTechniques{} techniques, split the NISS Cognitive Integrity metric into Cognitive
Reconnaissance and Cognitive Disruption (NISS~v1.1), normalized default weights, and added
neurorights mapping for 103 of \SnapshotTechniques{} techniques.
Version~1.6 (October 2026) replaces the hourglass stack with addressing version 2; adds the
evidence rule, the three scoring states, the retirement of the confirmed status, the hazards
dataset and NISS's stated limits; replaces case studies and representative techniques that
did not match the registry; reports which snapshot figures do and do not reproduce; adds
arguments for why neurosecurity and neuroethics are needed, with twenty newly cited
sources; qualifies or removes priority claims, universal negatives and diagnostic language;
removes six figures whose hard-coded values or labels no longer matched the tables or the
model; and condenses the future-work list.

\noindent\textbf{Note on AI use and preprint status:} I did not find anything
in Zenodo's terms of use that prohibits using AI to help write a preprint, so
I used it, heavily, to get this initial version out the door. The AI helped
me move faster, but the ideas, the architecture, the scoring, the clinical
mappings, those are mine. This preprint is me planting a seed. The real work
is what comes next: writing it properly with human collaborators, getting the
DSM-5-TR mappings reviewed by actual clinicians, running empirical validation,
and revising it until it holds up to scrutiny. There is a lot more to do.

\subsection{Future Work}

\begin{enumerate}
  \item \textbf{Re-examination of existing scores}: apply the evidence rule to the
        \LiveNissScored{} scored techniques
  \item \textbf{Expert review}: a neuroanatomist's review of the addressing model before
        the routes are integrated, and clinical review of the category associations
  \item \textbf{Interrater reliability and weight calibration}: a formal study with
        cybersecurity and neuroscience experts, with sensitivity analysis of the weights
  \item \textbf{Checks in the build}: run the band check, a score recomputation and the
        route gate in continuous integration
  \item \textbf{Standards input}: offer NISS to FIRST.org's CVSS community only after at
        least one independent researcher has reviewed the scoring method; no submission has
        been made
  \item \textbf{Empirical testing}: testing against BCI hardware in controlled environments
        with ethics approval
  \item \textbf{Exploratory directions} documented in the repository and not claims of this
        paper, each requiring ethics review before any implementation: hash-chained audit
        logs for interface firmware; an EEG dataset pipeline for validating components
        (\url{https://github.com/qinnovates/qinnovate/blob/main/shared/EEG-DATA-PIPELINE.md});
        a commodity-hardware vision-prosthesis security prototype
        (\url{https://github.com/qinnovates/qinnovate/blob/main/docs/research/lidar-bci-bridge-technical-spec.md});
        and consent-gated public-safety concepts
\end{enumerate}
